\documentclass[11pt, a4paper]{article}

\usepackage[T1]{fontenc}
\usepackage[utf8]{inputenc}
\usepackage{tgpagella}
\usepackage[spanish, es-noshorthands]{babel}
\usepackage{amsmath, amssymb, bm}
\usepackage{tabularx}
\usepackage{booktabs}
\usepackage[most]{tcolorbox}
\usepackage[margin=2.5cm]{geometry}
\usepackage{ragged2e}
\usepackage{fancyhdr}

\pagestyle{fancy}
\fancyhf{}
\setlength{\headheight}{16pt}
\lhead{\textbf{UCB Sede Tarija} | Ing. Mecatrónica}
\chead{}
\rhead{IMT-121 Circuitos Electrónicos I | Exit Ticket S10-02}
\lfoot{Prof: M.Sc. Ing. Bernardo Quiroga Turdera}
\rfoot{Página \thepage}
\renewcommand{\headrulewidth}{0.4pt}

\definecolor{ucbblue}{RGB}{0, 51, 102}

\begin{document}

\begin{tcolorbox}[colback=ucbblue!5!white, colframe=ucbblue, title=\textbf{Exit Ticket Formativo: Introducción a Circuitos Dinámicos (EC3)}]
    \centering \textbf{Cierre de Sesión (Sin Calificación Sumativa)}
\end{tcolorbox}

\noindent\textit{Nombre del Estudiante:} \hrulefill

\vspace{0.4cm}
\textbf{1.} Si el voltaje almacenado en un capacitor, instantes antes de accionar un interruptor es $v_C(0^-) = 5\,\text{V}$, ¿Cuál será su voltaje inmediatamente después de cerrar el circuito ($v_C(0^+)$)? \\
\rule{\textwidth}{0.4pt}

\vspace{0.6cm}
\textbf{2.} Si la corriente circulante por un inductor, previo a conmutar una fuente DC, es $i_L(0^-) = -100\,\text{mA}$, ¿Cuál será su corriente en el instante inmediatamente posterior ($i_L(0^+)$)? \\
\rule{\textwidth}{0.4pt}

\vspace{0.6cm}
\textbf{3.} Para un circuito RC serie simple compuesto por $R = 5\,\text{k}\Omega$ y un capacitor $C = 20\,\mu\text{F}$, calcule numéricamente su constante de tiempo $\tau$ indicando las unidades. Si decidimos cambiar el capacitor a uno de $100\,\mu\text{F}$, ¿El circuito responderá más rápido o más lento al transitorio? \\
\rule{\textwidth}{0.4pt} \\
\rule{\textwidth}{0.4pt}

\vspace{0.6cm}
\textbf{4.} Una variable transitoria $x(t)$ de un circuito arranca con una condición inicial $x(0^+) = 2\,\text{A}$ y posee un estado estacionario final hacia el infinito de $x(\infty) = 8\,\text{A}$. Describa cualitativamente el comportamiento de su curva a lo largo del tiempo (dónde empieza, hacia dónde va y cómo se mueve). \\
\rule{\textwidth}{0.4pt} \\
\rule{\textwidth}{0.4pt}

\end{document}
